\documentclass{article}

\usepackage{booktabs}
\usepackage{tabularx}
\usepackage{enumitem}

\title{Development Plan\\\progname}

\author{\authname}

\date{}

\input{../Comments.text}
\input{../Common.text}

\begin{document}

\maketitle

\begin{table}[hp]
\caption{Revision History} \label{TblRevisionHistory}
\begin{tabularx}{\textwidth}{llX}
\toprule
\textbf{Date} & \textbf{Developer(s)} & \textbf{Change}\\
\midrule
Apr 20, 2026 & Mirhoseininejad & Consolidated Development Plan documenting the process actually followed during the January--April 2026 CAS~741 project cycle for \progname{}.\\
\bottomrule
\end{tabularx}
\end{table}

\newpage{}

This Development Plan records the administrative, process, and technology
decisions adopted for \progname{}, a scientific computing project delivered
for CAS~741 between January and April 2026. The plan is authored as a solo
deliverable; sections that are framed in the course template around
multi-person teams are adapted here to reflect a single-developer workflow.
The remainder of the document covers: confidentiality and IP, copyright,
meeting and communication patterns with the supervisor, role allocation,
git workflow, project decomposition and scheduling, the Proof-of-Concept
plan, the expected technology stack, and the adopted coding standard.

\section{Confidential Information}

\progname{} contains no confidential information from industry. All source
code, documentation, and derived artefacts are developed from publicly
available libraries (MediaPipe Pose, MoveNet, Django, NumPy, SciPy) and from
publicly documented kinematics. No non-disclosure agreement is in place.

\section{IP to Protect}

There is no intellectual property requiring protection beyond the author's
own work. No third-party IP agreement applies, and the author is not
required to sign the Intellectual Property Guide Acknowledgement for this
project.

\section{Copyright License}

\progname{} is released under the MIT License, as recorded in the
\texttt{LICENSE} file at the root of the project repository
(\url{https://github.com/mmdmirh/cas741}). The MIT License was adopted
because it permits reuse of the refactored kinematics and state-machine
modules in future personal or academic projects without additional
licensing friction.

\section{Team Meeting Plan}

This is a solo project, so ``team meetings'' within the course sense do
not apply. The developer's working cadence was:

\begin{itemize}
  \item \textbf{Self-directed work sessions:} 4--6 focused sessions per
    week during term, primarily aligned with the CAS~741 Tuesday/Thursday
    lecture and tutorial slots, and weekends.
  \item \textbf{Supervisor / instructor interaction:} Dr.\ Spencer Smith
    (CAS~741 instructor) served as the primary reviewer. Interaction took
    the form of (i) in-class lecture Q\&A, (ii) the POC demonstration
    (2026-03-05), (iii) written peer-review issues opened on the project
    repository, and (iv) final-documentation review via GitHub issue \#27.
  \item \textbf{Peer-review interaction:} Structured bi-directional peer
    review with Yibing Mei, conducted on each major deliverable (SRS, MG,
    MIS, VnV Plan, VnV Report) via GitHub issues on the respective
    repositories. All feedback from the peer reviewer was triaged and
    either applied or explicitly noted as deferred.
\end{itemize}

\section{Team Communication Plan}

\begin{itemize}
  \item \textbf{Primary channel:} GitHub issues on
    \texttt{mmdmirh/cas741} --- used for all review requests, peer-review
    feedback, and instructor follow-ups. Each peer review was a separate
    issue to preserve traceability.
  \item \textbf{Deliverable hand-offs:} review-request issues were opened
    with a direct link to the relevant document PDF at its commit SHA.
  \item \textbf{Synchronous channel:} in-class discussion during lecture
    and tutorial slots; email for administrative matters.
\end{itemize}

\section{Team Member Roles}

\progname{} is developed by a single author, so all roles are filled by
Seyedmohamad Mirhosseininejad. For traceability the project was structured
around the following role-hats the author wore at different times:

\begin{itemize}
  \item \textbf{Requirements engineer} --- authored the SRS, scoped
    R1--R5 and NFR1--NFR2, maintained the traceability matrix.
  \item \textbf{Architect} --- authored the Module Guide (MG) and the
    Module Interface Specification (MIS); defined the eight-module
    decomposition (M1--M8).
  \item \textbf{Implementer} --- wrote the Python/Django backend, the
    JavaScript frontend, and the WebSocket transport layer.
  \item \textbf{Test engineer} --- authored the VnV Plan and the VnV
    Report; implemented \texttt{src/backend/tests/test\_vnv.py}
    (11 tests) and the interface-conformance harness.
  \item \textbf{Reviewer} --- provided peer-review feedback to Yibing
    Mei's project, and triaged incoming peer-review issues on this project.
\end{itemize}

\section{Workflow Plan}

\begin{itemize}
  \item \textbf{Git usage.} All work tracked in git on
    \texttt{github.com/mmdmirh/cas741}. The \texttt{main} branch holds the
    canonical state. For this solo project, small changes are committed
    directly to \texttt{main}; larger or more disruptive changes (such as
    the initial modular refactor from the CAS~772 prototype) are staged on
    short-lived topic branches and merged once the corresponding tests
    pass locally.
  \item \textbf{Pull requests.} Not used as a review gate for this solo
    project; peer review is conducted via issues rather than PRs, because
    each reviewer reads the published PDF rather than an open diff.
  \item \textbf{Issue management.} Every review request and every piece
    of peer-review feedback is filed as a GitHub issue with a descriptive
    title of the form \textit{``Peer Review of [artefact] by [reviewer]
    --- [topic]''}. Issues are closed with a comment citing the commit SHA
    that addressed them.
  \item \textbf{CI/CD.} A GitHub Actions workflow in \texttt{.github/}
    handles PDF rebuilds for the documentation site; the Python unit
    suite is run locally via \texttt{make test} before significant
    pushes. No production deployment pipeline is in scope for a
    research-grade project.
\end{itemize}

\section{Project Decomposition and Scheduling}

\begin{itemize}
  \item \textbf{GitHub Projects:} not used as a separate board; issue
    labels and milestones on the repository itself are sufficient for a
    solo workload.
  \item \textbf{Repository:}
    \url{https://github.com/mmdmirh/cas741}.
\end{itemize}

Scheduling followed the CAS~741 course outline deadlines. The big-picture
schedule, as actually delivered, was:

\begin{table}[h!]
\centering
\begin{tabularx}{\textwidth}{l l X}
\toprule
\textbf{Phase} & \textbf{Target (2026)} & \textbf{Artefact}\\
\midrule
Problem Statement          & Jan                & ProblemStatement.pdf\\
SRS (draft + peer review)  & Feb 11--15         & SRS.pdf\\
VnV Plan                   & Feb 18 -- Mar 12   & VnVPlan.pdf\\
Proof-of-Concept demo      & Mar 5              & live demo\\
MG + MIS                   & Apr 3--5           & MG.pdf, MIS.pdf\\
VnV Report                 & Apr 8--15          & VnVReport.pdf\\
Reflection and Trace       & Apr                & ReflectAndTrace.pdf\\
Final documentation pass   & Apr 15--20         & All PDFs\\
EXPO presentation          & Apr                & expo.pdf\\
\bottomrule
\end{tabularx}
\caption{Project Schedule, as Delivered}
\label{tab:schedule}
\end{table}

\section{Proof of Concept Demonstration Plan}

The two principal project risks identified at the start of the term were:

\begin{enumerate}
  \item \textbf{Pose-estimation feasibility risk.} Could a commodity
    webcam and a third-party pose-estimation library (MediaPipe or
    MoveNet) deliver skeletal keypoints at a stable enough frame rate,
    and with enough angular accuracy, to support real-time exercise
    form-feedback?
  \item \textbf{Architectural-feasibility risk.} Could the monolithic
    CAS~772 prototype be re-expressed as a hidden-secret modular
    decomposition (eight modules M1--M8, with abstract-base-class
    interfaces) without losing real-time behaviour?
\end{enumerate}

The POC demonstration (2026-03-05) addressed both risks by showing
end-to-end live operation: a browser client capturing webcam frames,
the Django WebSocket backend streaming MediaPipe landmarks back,
joint angles computed via the new \texttt{KinematicEngine}, and the
\texttt{ExerciseStateMachine} counting valid repetitions with visible
AR overlay on the frontend. Surviving the POC confirmed both risks
were retired and permitted the project to move into the design-artefact
phase (MG / MIS / VnV Report).

\section{Expected Technology}

\subsection*{Programming languages}
\begin{itemize}
  \item \textbf{Backend:} Python 3.9.6.
  \item \textbf{Frontend:} standard HTML / CSS / JavaScript (ES2020);
    no transpilation step, deliberately kept dependency-free to preserve
    the NFR2 portability claim.
\end{itemize}

\subsection*{Frameworks and libraries}
\begin{itemize}
  \item \textbf{Django + Django Channels} --- WebSocket transport between
    browser client and backend pipeline.
  \item \textbf{MediaPipe Pose} --- production pose-estimation backend
    (M4). \textbf{MoveNet} --- comparison backend used in the offline
    dynamic benchmark.
  \item \textbf{NumPy, SciPy} --- vector arithmetic for
    \texttt{KinematicEngine.compute\_angle} and for the Kalman smoother
    in M8.
\end{itemize}

\subsection*{Pre-trained models}
MediaPipe Pose and MoveNet are used as provided by their upstream
distributions. No custom model training is in scope.

\subsection*{Components implemented despite library availability}
The repetition state machine (M6) is hand-implemented rather than
adapted from a general-purpose FSM library, because the validity
criteria (IM\_SquatValid, IM\_BicepCurlValid, $t_{min}$/$t_{max}$
timing guards) are tightly coupled to the SRS instance models and
benefit from being readable alongside the requirement text.

\subsection*{Linter}
\texttt{flake8} and \texttt{pylint} are run locally on the Python
backend as part of the static verification workflow recorded in the
VnV Plan.

\subsection*{Unit testing framework}
\texttt{pytest 8.4.2}. Tests live in
\texttt{src/backend/tests/test\_vnv.py} and an interface-conformance
harness in \texttt{src/backend/test\_interfaces.py}; both are invoked
by \texttt{make test}.

\subsection*{Code coverage tool}
\texttt{coverage.py} (pytest-cov plug-in) was investigated for the
Code Coverage Metrics section of the VnV Report.

\subsection*{Continuous integration}
A GitHub Actions workflow (\texttt{.github/workflows}) rebuilds the
documentation site on push. The Python test suite is run locally before
pushes; a production CI pipeline was not in scope for a research-grade,
solo project.

\subsection*{Performance measurement}
Average-latency and frame-to-frame stability measurements were
produced by a bespoke offline benchmark script rather than by a
general-purpose profiler such as Valgrind, because the metrics of
interest (per-backend frame latency and $\overline{|\Delta\theta|}$)
are domain-specific.

\subsection*{Other tools}
git, GitHub, GitHub Issues, and GitHub Pages (for hosted PDF review
links) are part of the development toolchain.

\section{Coding Standard}

\begin{itemize}
  \item \textbf{Python:} PEP~8, enforced by \texttt{flake8}. Type hints
    are added on public module interfaces (M3, M4, M5, M6, M8) to make
    the MIS contracts visible at the language level, backed by
    \texttt{abc.ABC} base classes.
  \item \textbf{JavaScript:} vanilla ES2020 with single-responsibility
    modules (\texttt{m2\_display\_output.js},
    \texttt{m7\_ui\_rendering.js}) and no framework dependency.
  \item \textbf{LaTeX:} standard article class across all CAS~741
    deliverables; shared macros for \texttt{\textbackslash progname} and
    \texttt{\textbackslash authname} are factored into
    \texttt{docs/Common.text} and \texttt{docs/Comments.text}.
\end{itemize}

\end{document}
